% !TeX program = lualatex
% Definition and sources only; this file is not an LLM solution attempt.
\documentclass[11pt,a4paper]{article}
\usepackage[margin=25mm]{geometry}
\usepackage{fontspec}
\setmainfont{lmroman10-regular.otf}[BoldFont=lmroman10-bold.otf,
 ItalicFont=lmroman10-italic.otf,BoldItalicFont=lmroman10-bolditalic.otf]
\usepackage{amsmath,amssymb,mathtools}
\usepackage{unicode-math}
\setmathfont{latinmodern-math.otf}
\AtBeginDocument{\let\setminus\smallsetminus}
\usepackage{xurl,hyperref}
\hypersetup{colorlinks=true,urlcolor=blue}
\setcounter{secnumdepth}{0}
\setlength{\parindent}{0pt}
\setlength{\parskip}{5pt}
\setlength{\emergencystretch}{3em}
\begin{document}
\section*{Exact recursion and local profile stability for the broken tetrahedron construction}\label{statement}

\subsection{Definitions and mathematical statement}
Let $F=K_4^-$ be the simple $3$-uniform hypergraph on $\{1,2,3,4\}$
with edges $123$, $124$, and $134$, where, for example, $123$ means
$\{1,2,3\}$. Thus $F$ is a complete $3$-uniform hypergraph on four
vertices with one triple deleted.
A $3$-uniform hypergraph $H$ is $F$-free if it has no injective copy
of these three edges. The missing fourth triple is unrestricted in
a copy, so four vertices spanning all four triples also violate
$F$-freeness. Equivalently, every set of four vertices of an $F$-free
hypergraph spans at most two edges.

For $n\ge4$, set
\[
 \operatorname{ex}(n,F)=\max\{|E(H)|:|V(H)|=n, H\text{ is }F\text{-free}\},
 \qquad
 \pi(F)=\lim_{n\to\infty}\frac{\operatorname{ex}(n,F)}{\binom n3}.
\]
The maximum ranges over simple hypergraphs; repeated edges are not
counted. The normalized extremal densities have a limit by averaging
over vertex subsets.
The workshop conjecture is $\pi(F)=2/7$. Accordingly, a resolution
requires the sharp asymptotic extremal density, not an exact formula
for every finite $n$. The proposed upper bound says that for every
$\varepsilon>0$, all sufficiently large triple systems of density
exceeding $2/7+\varepsilon$ have four vertices spanning at least three
triples. The lower side asks for $F$-free systems approaching density
$2/7$.
\subsection{Short English statement}
Is two sevenths the Turán density of a tetrahedron with one triple removed?
\subsection{Sources}
\begin{itemize}
\item[S1] ULAM.AI and the UnsolvedMath Contributors, supplied problems.json, catalogue entry, AIM Workshop Problem Lists. Catalogue identity and source transcription; CC BY 4.0.
\url{https://huggingface.co/datasets/ulamai/UnsolvedMath}.
\item[S2] American Institute of Mathematics, Hypergraph Turan problem, item 11.4. Original workshop question.
\url{http://aimpl.org/hypergraphturan/1/}.
\end{itemize}
\end{document}
